% NOTE (2026-09-25): superseded by ../draft.tex, which was revised after a 46-item review; this file is kept as a source only.
% =====================================================================================
% sections/intro.tex -- title, abstract and Introduction of the short draft of "Degrees of Separation".
% Author: Sora Yng. Written 2026-09-25. Meant to be \input right after \begin{document}; it sets the title and
% calls \maketitle, so the including file should not call \maketitle again.
% Numbers: every number is copied from paper/main.tex (Introduction; Results, sec:map to sec:boundary;
% Table tab:ladder); none is new.
%   +0.419 [+0.355, +0.482], 52 fields ........ main.tex, "The two orderings agree" (sec:map)
%   +0.43 -> +0.14 ............................ main.tex, control ladder (sec:status; Table tab:ladder, 44 fields)
%   0.84% per SD, interval 0 to 1.6%, ~3% ..... main.tex, within institutions (+0.0084 [+0.0000, +0.0164]) and
%                                               reliability-corrected remainder (up to about 3%)
%   66% [51, 79], 28 UK subjects .............. main.tex, UK prior-attainment control (sec:uk)
%   +0.031 [+0.021, +0.041]; +0.026 .......... main.tex, career time (sec:career; 39 US fields, 33 UK subjects)
%   three fifths shared with selectivity ...... main.tex, sec:career (0.39 of the raw slope remains)
%   +0.120 [-0.033, +0.258] vs +0.523 ......... main.tex, sec:boundary
%   about a fifth between employer sectors .... main.tex, T8 (0.21 [0.08, 0.34])
% Citations: keys in ../refs_merged.bib, plus anelli2020returns in refs_short.bib (checked again against
% api.crossref.org/works/10.1093/jeea/jvz070 on 2026-09-25: title, author, journal, 18(6), 2824-2868 match).
% The theory section is referred to as Section~\ref{sec:theory}.
% =====================================================================================
\providecommand{\ci}[2]{$[#1,\allowbreak\,#2]$}

\title{Academic and employer reputation agree at the university level, much less at the department}
\author{Sora Yng}
\date{}
\maketitle

\begin{abstract}
When students choose a university and a field, which reputation do they buy? A university $\times$ field programme
is valued in two markets; we read its academic reputation from faculty hiring and its employer reputation from
graduates' pay. Using public data on 52 US fields and 33 UK subjects, we ask where the two rankings agree. Within
fields they agree, but the agreement sits at the university level. Most of it is shared with intake selectivity:
controls for selectivity and institutional traits lower mean US agreement from $+0.43$ to $+0.14$. Within a
university, one standard deviation of a department's field rank goes with about 0.8\% higher pay, an interval
reaching zero; only computer science is clear on its own. Agreement rises with years since graduation and is weak
where pay follows national scales. The design is descriptive; for students, the reputation both markets recognise is
mostly the university's.
\end{abstract}

\section{Introduction}\label{sec:intro}

A student who applies to a university and a field chooses a programme: economics at one university, nursing at
another. Much of what that choice promises is reputation, and reputation is valued by others. The returns to such
choices are now estimated programme by programme. \citet{anelli2020returns} compares applicants just above and just
below the admission cutoff of a highly selective private university, and \citet{kirkeboen2016field} use the cutoffs
that ration entry to fields. These designs estimate what admission is worth, not which level of the programme
carries the reputation that others value: the university's name or the department.

A programme's reputation in a market is that market's expectation of the quality of what it buys from the programme
\citep{tadelis1999name, macleod2015reputation}. A university $\times$ field programme sells in two markets under one
name. In the academic labour market, hiring committees buy its doctoral graduates as faculty. The net direction of
their hires orders departments in a steep hierarchy \citep{clauset2015systematic}. We read that ordering as academic
reputation, measured by the field rank $F$ and university-wide rank $G$ of \citet{wapman2022quantifying}. In the
graduate labour market, employers buy its bachelor's graduates. What they pay reveals employer reputation, mixed
with what they later learn about graduates' own productivity and with who enrols; we measure it by programme median
earnings. Students, a third party, reveal a third ranking through their admission choices: intake selectivity.

Research on the returns to college has mostly measured standing by selectivity, rankings or brand. Among students
admitted to similar colleges, the payoff to the more selective one is small on average \citep{dale2002estimating,
dale2014estimating}. The causal payoff to elite colleges concentrates in top jobs and incomes
\citep{zimmerman2019elite, chetty2023diversifying}. The pay gradient in selectivity differs by field, nearly flat in
nursing and teacher education and steep in computer science and economics \citep{bloem2024understanding}. In the
United Kingdom, net of intake attainment, no subject's league-table rating correlates positively and significantly
with returns \citep{britton2022how}. Whether the academy's own ranking of a department agrees with employers'
valuation of its graduates, and at which level, has not to our knowledge been measured.

We measure that agreement as coupling: the Spearman correlation, across the universities that teach a field, between
a department's field rank and its graduates' median earnings. We map it for 52 US fields on College Scorecard
earnings \citep{scorecard}, with Census PSEO as a second source \citep{pseo, pseo2026}. For 33 UK subjects we set a
university-wide hiring rank from ORCID-derived career records \citep{li2026orcid} against LEO earnings, which are
also released by graduates' prior attainment \citep{dfe_leo_provider}. All inputs are public, and the analysis is
descriptive. A short model (Section~\ref{sec:theory}) states when hiring reveals academic reputation and pay reveals
employer reputation, and what coupling then measures. We wrote it after most results were known. Its only content
not fitted to them is nine tests, T0--T8, specified before the analysis; not publicly registered.

Within fields the two rankings agree: across 52 US fields mean coupling is $+0.419$ \ci{+0.355}{+0.482}. The
agreement sits at the university level, and most of it is shared with intake selectivity. Controls for selectivity,
Pell share, public or private control and state earnings level lower mean coupling from $+0.43$ to $+0.14$. The
field rank does not detectably out-predict the university-wide rank. Within universities, one standard deviation of
field rank goes with about 0.8\% higher four-year pay (95\% interval 0 to 1.6\%). Corrected for error in the ranks
it is up to about 3\%, so the department part is diluted, not shown absent. It is clear on its own only in computer
science. In the UK data, which have no department rank, 66\% \ci{51}{79} of coupling remains among graduates of the
same prior-attainment band.

Within fixed US graduation cohorts, coupling rises by $+0.031$ per year since graduation \ci{+0.021}{+0.041}, and UK
coupling rises by about as much. The rise mixes career time with calendar time, and about three fifths of the US
rise is shared with selectivity. Agreement is weak where pay follows national scales: UK nursing, medicine and
teaching couple at $+0.120$ \ci{-0.033}{+0.258}, against $+0.523$ for the other subjects. Where a schedule sets pay,
weak coupling is not evidence that employers rank universities differently.

Nothing here shows that attending a programme raises pay, since the design does not separate value added, selection
and employer screening. The field rank is academic reputation only under the model's assumptions, and an audit finds
that it partly reflects doctoral production. On the pay side, about a fifth of coupling runs between employer
sectors, which reflects employers' valuation only if they ration access to sectors. Nor do we identify why agreement
below the university is weak. Departments may carry little that employers price, or employers may attribute little
to departments; no test here separates the two.

For research on the returns to college, the results locate the reputation at stake. Admission to a programme confers
a university's name and a department's standing. The academy and employers rank the name alike, largely through who
enrols, and the department's standing only faintly. Designs that compare attending more and less selective
universities \citep{dale2002estimating, anelli2020returns} therefore act on the margin where the two markets agree.
For a student, a department's field rank says little about graduates' pay beyond the university's rank and intake,
except perhaps in computer science. Whether a higher-ranked university or department raises pay is a causal question
these data cannot answer. It needs admission-based variation \citep{hoekstra2009effect, kirkeboen2016field,
anelli2020returns}; this paper supplies the descriptive map such work can start from.
